\documentclass[conference,twocolumn]{IEEEtran}

% Packages
\usepackage[utf8]{inputenc}
\usepackage[T1]{fontenc}
\usepackage{kotex}
\usepackage{amsmath,amssymb,amsfonts}
\usepackage{graphicx}
\usepackage{textcomp}
\usepackage{xcolor}
\usepackage{booktabs}
\usepackage{algorithm}
\usepackage{algorithmic}
\usepackage{multirow}
\usepackage{cite}
\usepackage{hyperref}
\usepackage{array}
\usepackage{subfig}

% Define colors
\definecolor{codegreen}{rgb}{0,0.6,0}
\definecolor{codegray}{rgb}{0.5,0.5,0.5}
\definecolor{codepurple}{rgb}{0.58,0,0.82}

\begin{document}

\title{Bearing Remaining Useful Life Prediction using CNN-LSTM with Advanced Signal Processing}

\author{
\IEEEauthorblockN{Jinwook Kim}
\IEEEauthorblockA{Department of Mechanical Engineering\\
Seoul National University of Science and Technology\\
Seoul, Republic of Korea\\
Email: jinwook@seoultech.ac.kr}
\and
\IEEEauthorblockN{Seungil Seo}
\IEEEauthorblockA{Department of Mechanical System Design Engineering\\
Seoul National University of Science and Technology\\
Seoul, Republic of Korea\\
Email: seungil@seoultech.ac.kr}
}

\maketitle

\begin{abstract}
Accurate prediction of remaining useful life (RUL) for rolling element bearings is critical for implementing effective predictive maintenance strategies in industrial machinery. This paper presents a hybrid deep learning approach combining Convolutional Neural Networks (CNN) and Long Short-Term Memory (LSTM) networks for bearing RUL prediction, integrated with advanced signal processing techniques. Our methodology employs band-pass filtering (1000-5000 Hz), Hilbert transform-based envelope analysis, and Daubechies-4 (Db4) wavelet decomposition to extract meaningful degradation features from raw vibration signals. The extracted features include Root Mean Square (RMS), entropy measures, and wavelet coefficients (D4\_RMS, D5\_RMS, D5\_Entropy) that effectively capture the bearing degradation progression. The proposed CNN-LSTM architecture learns both spatial and temporal patterns from multi-channel vibration data sampled at 25.6 kHz. We validated our approach on the KSPHM-KIMM 2025 Bearing RUL Prediction Challenge dataset, achieving 7th place among 70 participating teams. Experimental results demonstrate that the combination of physics-informed signal processing with deep learning enables robust RUL prediction even under varying operational conditions.
\end{abstract}

\begin{IEEEkeywords}
Bearing fault diagnosis, Remaining useful life, CNN-LSTM, Wavelet transform, Envelope analysis, Predictive maintenance
\end{IEEEkeywords}

\section{Introduction}

Rolling element bearings are among the most critical components in rotating machinery, and their unexpected failures can lead to significant economic losses, safety hazards, and unplanned downtime in industrial systems \cite{lei2018machinery}. According to industrial surveys, bearing failures account for approximately 40-50\% of all rotating machinery failures \cite{randall2011rolling}. Therefore, developing accurate remaining useful life (RUL) prediction methods for bearings has become a crucial research topic in the field of prognostics and health management (PHM).

Traditional bearing condition monitoring approaches rely on threshold-based methods using statistical features extracted from vibration signals \cite{jardine2006review}. While these methods are interpretable and computationally efficient, they often fail to capture the complex nonlinear degradation patterns exhibited by bearings under realistic operating conditions. Furthermore, the degradation process of bearings is influenced by multiple factors including load variations, speed changes, and environmental conditions, making accurate RUL prediction challenging.

Recent advances in deep learning have opened new possibilities for bearing prognosis. Convolutional Neural Networks (CNNs) have demonstrated remarkable capability in automatically extracting hierarchical features from raw or minimally processed signals \cite{zhang2017new}. Long Short-Term Memory (LSTM) networks excel at capturing long-term temporal dependencies in sequential data, making them well-suited for modeling degradation trajectories \cite{zhao2017machine}. The combination of these architectures, known as CNN-LSTM, leverages the strengths of both approaches: CNN layers extract spatial features from signal segments while LSTM layers model the temporal evolution of these features over the bearing's lifetime.

However, applying deep learning directly to raw vibration signals may not yield optimal results due to the presence of noise and irrelevant frequency components. Signal processing techniques such as band-pass filtering, envelope analysis, and wavelet transforms can enhance the signal-to-noise ratio and extract features that are more sensitive to bearing degradation \cite{antoni2007fast}. The integration of domain knowledge through physics-informed signal processing with data-driven deep learning models represents a promising direction for improving RUL prediction accuracy.

In this paper, we present a comprehensive approach for bearing RUL prediction that combines advanced signal processing techniques with a CNN-LSTM deep learning architecture. Our main contributions are as follows:

\begin{itemize}
    \item We develop a signal processing pipeline incorporating band-pass filtering (1000-5000 Hz), Hilbert transform-based envelope analysis, and Db4 wavelet decomposition optimized for bearing fault feature extraction.
    \item We extract physics-informed features including RMS and entropy from multiple wavelet scales (D4, D5) that effectively capture bearing degradation progression.
    \item We design a CNN-LSTM architecture that learns both spatial and temporal patterns from multi-channel vibration data.
    \item We validate our approach on the KSPHM-KIMM 2025 Bearing RUL Prediction Challenge, achieving 7th place among 70 participating teams.
\end{itemize}

The remainder of this paper is organized as follows. Section II reviews related work on bearing fault diagnosis and RUL prediction. Section III describes our methodology including signal processing techniques, feature extraction, and the CNN-LSTM model architecture. Section IV presents experimental setup and results. Section V concludes the paper with a discussion of findings and future research directions.

\section{Related Work}

\subsection{Bearing Fault Diagnosis and Prognosis}

Bearing fault diagnosis has been extensively studied using vibration analysis techniques. Traditional approaches focus on extracting time-domain statistical features (mean, variance, RMS, kurtosis, crest factor) and frequency-domain features (spectral energy, dominant frequencies) \cite{randall2011rolling}. The envelope analysis technique, based on the Hilbert transform, has been particularly effective for detecting periodic impulses generated by bearing defects \cite{mccormick1998cyclostationarity}.

Wavelet transform provides multi-resolution analysis capability, enabling the decomposition of vibration signals into different frequency bands. The Discrete Wavelet Transform (DWT) is widely used for bearing fault diagnosis due to its ability to capture transient events associated with bearing defects \cite{kumar2013wavelet}. Among various wavelet families, Daubechies wavelets, particularly Db4, have been shown to be effective for bearing vibration analysis \cite{toma2020bearing, nizwan2016wavelet}.

\subsection{Deep Learning for RUL Prediction}

Deep learning methods have gained significant attention for bearing RUL prediction in recent years. Convolutional Neural Networks have been applied to learn features directly from raw vibration signals or time-frequency representations \cite{li2018remaining}. LSTM networks have been employed to model the temporal evolution of bearing health indicators \cite{zhao2017machine}.

Hybrid architectures combining CNN and LSTM have shown promising results. The CNN layers serve as automatic feature extractors, while LSTM layers capture the temporal dependencies in the extracted features \cite{chen2020bearing}. Attention mechanisms have also been incorporated to focus on the most relevant time steps for RUL prediction \cite{wang2020remaining}.

Despite these advances, many deep learning approaches treat the problem as purely data-driven, potentially missing valuable insights from domain knowledge. Recent work has emphasized the importance of integrating physics-informed features with deep learning models to improve prediction accuracy and generalization \cite{li2020deep}.

\section{Methodology}

\subsection{Data Description}

The KSPHM-KIMM 2025 Bearing RUL Prediction Challenge provides run-to-failure vibration data from rolling element bearings. The dataset consists of 8 training sets and 6 validation sets, each representing a complete degradation process from healthy state to failure.

Each data sample is stored in TDMS format and contains:
\begin{itemize}
    \item Four-channel vibration signals (CH1-CH4) sampled at 25,600 Hz
    \item Operating condition data (torque, temperature) sampled at lower rates
    \item Each acquisition period captures 10 seconds of data at 10-minute intervals
\end{itemize}

The vibration data contains 256,000 data points per channel per file (25,600 Hz $\times$ 10 seconds). The multi-channel configuration enables comprehensive monitoring of bearing vibration from different sensor positions.

\subsection{Signal Processing}

Our signal processing pipeline consists of three main stages: band-pass filtering, envelope analysis, and wavelet decomposition. These techniques are specifically designed to enhance bearing fault-related features while suppressing noise.

\subsubsection{Band-Pass Filter}

Low-frequency components (below 1000 Hz) in vibration signals are often dominated by structural vibrations and external disturbances, while high-frequency components (above 5000 Hz) are susceptible to noise. Based on the analysis of bearing fault frequencies and their harmonics, we apply a 4th-order Butterworth band-pass filter with cutoff frequencies of 1000 Hz and 5000 Hz.

The transfer function of an $n$-th order Butterworth band-pass filter is given by:
\begin{equation}
    H(s) = \frac{H_0 \cdot (\omega_0/Q)s}{s^2 + (\omega_0/Q)s + \omega_0^2}
\end{equation}
where $\omega_0 = \sqrt{\omega_L \cdot \omega_H}$ is the center frequency, $Q = \omega_0/(\omega_H - \omega_L)$ is the quality factor, and $\omega_L$, $\omega_H$ are the lower and upper cutoff frequencies, respectively.

The Butterworth filter is characterized by maximally flat magnitude response in the passband:
\begin{equation}
    |H(j\omega)|^2 = \frac{1}{1 + \left(\frac{\omega}{\omega_c}\right)^{2n}}
\end{equation}
where $n$ is the filter order and $\omega_c$ is the cutoff frequency.

\subsubsection{Envelope Analysis}

Envelope analysis based on the Hilbert transform is employed to extract the amplitude modulation pattern of vibration signals. When a bearing defect occurs, periodic impulses are generated as rolling elements pass over the defect. These impulses excite the resonance frequencies of the bearing and housing, resulting in amplitude-modulated signals.

The Hilbert transform of a signal $x(t)$ is defined as:
\begin{equation}
    \hat{x}(t) = \mathcal{H}\{x(t)\} = \frac{1}{\pi} \int_{-\infty}^{\infty} \frac{x(\tau)}{t - \tau} d\tau
\end{equation}

The analytic signal is constructed as:
\begin{equation}
    z(t) = x(t) + j\hat{x}(t)
\end{equation}

The envelope (instantaneous amplitude) is obtained as:
\begin{equation}
    A(t) = |z(t)| = \sqrt{x^2(t) + \hat{x}^2(t)}
\end{equation}

This envelope signal captures the modulation pattern caused by bearing defects, enabling the detection of fault characteristic frequencies even when the signal is contaminated by noise.

\subsubsection{Wavelet Transform}

The Discrete Wavelet Transform (DWT) provides multi-resolution analysis of the vibration signal, decomposing it into approximation and detail coefficients at different scales. We employ the Daubechies-4 (Db4) wavelet, which has been shown to be effective for bearing vibration analysis due to its good balance between time and frequency localization \cite{kumar2013wavelet}.

The DWT of a signal $x(t)$ is computed as:
\begin{equation}
    W_{\psi}(j, k) = \frac{1}{\sqrt{2^j}} \int_{-\infty}^{\infty} x(t) \psi^*\left(\frac{t - 2^j k}{2^j}\right) dt
\end{equation}
where $\psi(t)$ is the mother wavelet, $j$ is the scale parameter, and $k$ is the translation parameter.

At decomposition level 5, we obtain detail coefficients D1 through D5, corresponding to different frequency bands. Based on our analysis, we focus on D4 and D5 coefficients as they capture the frequency ranges most relevant to bearing fault characteristics.

\subsection{Feature Extraction}

From the processed signals, we extract features that effectively characterize bearing degradation:

\subsubsection{Root Mean Square (RMS)}

RMS is a fundamental indicator of signal energy and vibration amplitude:
\begin{equation}
    RMS = \sqrt{\frac{1}{N} \sum_{i=1}^{N} x_i^2}
\end{equation}
where $N$ is the number of samples and $x_i$ represents the signal values.

RMS is computed from:
\begin{itemize}
    \item Band-pass filtered signal (CH2\_1000-5000Hz\_BPF)
    \item Envelope of band-pass filtered signal (CH2\_1000-5000Hz\_RMS\_Envelope)
    \item Wavelet detail coefficients (D4\_RMS, D5\_RMS)
\end{itemize}

\subsubsection{Shannon Entropy}

Shannon entropy quantifies the complexity or irregularity of a signal:
\begin{equation}
    H = -\sum_{i} p(x_i) \log_2 p(x_i)
\end{equation}
where $p(x_i)$ is the probability density of the signal values.

Entropy increases as bearing degradation progresses due to the increasing irregularity and complexity of the vibration signal. We compute D5\_Entropy from the wavelet D5 coefficients using histogram-based probability estimation with 64 bins.

\subsubsection{Feature Summary}

The final feature set for each sample consists of:
\begin{itemize}
    \item CH2\_1000-5000Hz\_BPF: RMS of band-pass filtered CH2 signal
    \item CH2\_1000-5000Hz\_RMS\_Envelope: RMS of envelope signal
    \item D4\_RMS: RMS of D4 wavelet coefficients (all channels)
    \item D5\_RMS: RMS of D5 wavelet coefficients (all channels)
    \item D5\_Entropy: Entropy of D5 wavelet coefficients (all channels)
\end{itemize}

\subsection{Model Architecture}

We propose a CNN-LSTM architecture that learns both spatial features from multi-channel vibration data and temporal patterns from the degradation sequence. The architecture consists of three main components: CNN feature extractor, LSTM sequence encoder, and fully connected predictor.

\subsubsection{CNN Feature Extractor}

The CNN component consists of multiple convolutional layers with the following structure:
\begin{itemize}
    \item Conv1D layer with 64 filters, kernel size 3, ReLU activation
    \item Batch normalization and max pooling
    \item Conv1D layer with 128 filters, kernel size 3, ReLU activation
    \item Batch normalization and max pooling
    \item Conv1D layer with 256 filters, kernel size 3, ReLU activation
\end{itemize}

\subsubsection{LSTM Sequence Encoder}

The LSTM component captures temporal dependencies in the feature sequence. An LSTM cell computes:
\begin{equation}
    f_t = \sigma(W_f \cdot [h_{t-1}, x_t] + b_f)
\end{equation}
\begin{equation}
    i_t = \sigma(W_i \cdot [h_{t-1}, x_t] + b_i)
\end{equation}
\begin{equation}
    \tilde{C}_t = \tanh(W_C \cdot [h_{t-1}, x_t] + b_C)
\end{equation}
\begin{equation}
    C_t = f_t \odot C_{t-1} + i_t \odot \tilde{C}_t
\end{equation}
\begin{equation}
    o_t = \sigma(W_o \cdot [h_{t-1}, x_t] + b_o)
\end{equation}
\begin{equation}
    h_t = o_t \odot \tanh(C_t)
\end{equation}
where $f_t$, $i_t$, $o_t$ are the forget, input, and output gates, $C_t$ is the cell state, $h_t$ is the hidden state, and $\sigma$ denotes the sigmoid function.

Our architecture uses a bidirectional LSTM with 128 hidden units per direction, followed by a second LSTM layer with 64 units.

\subsubsection{Fully Connected Predictor}

The output from the LSTM is passed through fully connected layers:
\begin{itemize}
    \item Dense layer with 64 units, ReLU activation, dropout (0.3)
    \item Dense layer with 32 units, ReLU activation
    \item Output layer with 1 unit (RUL prediction)
\end{itemize}

\subsubsection{Loss Function}

We employ an asymmetric Mean Squared Error (MSE) loss function that penalizes late predictions more heavily than early predictions:
\begin{equation}
    \mathcal{L} = \frac{1}{N} \sum_{i=1}^{N}
    \begin{cases}
        \alpha (y_i - \hat{y}_i)^2, & \text{if } \hat{y}_i > y_i \\
        (y_i - \hat{y}_i)^2, & \text{otherwise}
    \end{cases}
\end{equation}
where $y_i$ is the true RUL, $\hat{y}_i$ is the predicted RUL, and $\alpha > 1$ is the penalty factor for late predictions. In practice, predicting failure later than actual occurrence is more dangerous than early prediction, justifying this asymmetric penalty.

\section{Experiments}

\subsection{Experimental Setup}

The experiments were conducted on a system with the following specifications:
\begin{itemize}
    \item Operating System: Linux (Ubuntu 20.04)
    \item Python: 3.8.10
    \item Deep Learning Framework: PyTorch 1.9.0
    \item GPU: NVIDIA RTX 3090 (24GB VRAM)
\end{itemize}

The dataset was split as follows:
\begin{itemize}
    \item Training: Train1-Train8 (8 run-to-failure experiments)
    \item Validation: Validation1-Validation6 (6 experiments for challenge evaluation)
\end{itemize}

\subsection{Training Details}

The model was trained with the following hyperparameters:
\begin{itemize}
    \item Optimizer: AdamW with weight decay 0.01
    \item Learning rate: 0.001 with cosine annealing scheduler
    \item Batch size: 32
    \item Number of epochs: 100
    \item Early stopping patience: 15 epochs
    \item Sequence length: 50 time steps (500 seconds of operation)
\end{itemize}

Data augmentation techniques included:
\begin{itemize}
    \item Random noise injection (SNR 40dB)
    \item Time shifting ($\pm$5\% of sequence length)
    \item Feature scaling using min-max normalization
\end{itemize}

The RUL labels were normalized to the range [0, 1] by dividing by the maximum RUL in each run-to-failure experiment. This normalization helps the model learn consistent degradation patterns across experiments with different lifetimes.

\subsection{Results}

\begin{table}[t]
\centering
\caption{Performance Comparison on Validation Set}
\label{tab:results}
\begin{tabular}{lcc}
\toprule
\textbf{Method} & \textbf{RMSE} & \textbf{Score} \\
\midrule
Baseline (RMS only) & 0.187 & 0.721 \\
CNN only & 0.142 & 0.803 \\
LSTM only & 0.128 & 0.835 \\
CNN-LSTM (raw signal) & 0.098 & 0.891 \\
\textbf{CNN-LSTM (proposed)} & \textbf{0.074} & \textbf{0.923} \\
\bottomrule
\end{tabular}
\end{table}

Table~\ref{tab:results} presents the performance comparison of different methods on the validation set. Our proposed CNN-LSTM model with advanced signal processing achieves the best performance with an RMSE of 0.074 and a score of 0.923.

The results demonstrate that:
\begin{enumerate}
    \item The combination of CNN and LSTM outperforms individual architectures
    \item Signal processing significantly improves prediction accuracy compared to using raw signals
    \item The proposed features (wavelet RMS and entropy) effectively capture degradation patterns
\end{enumerate}

\begin{table}[t]
\centering
\caption{Ablation Study on Feature Contributions}
\label{tab:ablation}
\begin{tabular}{lc}
\toprule
\textbf{Feature Configuration} & \textbf{RMSE} \\
\midrule
All features & 0.074 \\
w/o D5\_Entropy & 0.082 \\
w/o D4\_RMS & 0.086 \\
w/o Envelope RMS & 0.091 \\
w/o Band-pass RMS & 0.095 \\
Only raw channel RMS & 0.142 \\
\bottomrule
\end{tabular}
\end{table}

Table~\ref{tab:ablation} shows the ablation study results on feature contributions. All proposed features contribute to the final performance, with band-pass filtered RMS and envelope RMS being the most important features.

In the KSPHM-KIMM 2025 Bearing RUL Prediction Challenge, our approach achieved \textbf{7th place among 70 participating teams}, demonstrating the effectiveness of combining physics-informed signal processing with deep learning for bearing prognosis.

\subsection{Discussion}

Our experimental analysis revealed several important findings:

\textbf{Importance of CH2 channel:} Among the four vibration channels, CH2 consistently showed the clearest degradation trend across all training sets. The 1000-5000 Hz band-pass filtered signal from CH2 exhibited progressive RMS increase as bearings approached failure.

\textbf{Wavelet scale selection:} The D4 and D5 scales of Db4 wavelet decomposition proved most informative for RUL prediction. These scales correspond to frequency ranges where bearing fault harmonics are prominent, while maintaining good signal-to-noise ratio compared to higher-frequency detail coefficients.

\textbf{Envelope analysis effectiveness:} The Hilbert transform-based envelope analysis successfully extracted the amplitude modulation pattern caused by bearing defects, even when the original signal was contaminated by noise. This preprocessing significantly enhanced the discriminability of health indicators.

\textbf{Asymmetric loss function:} The asymmetric MSE loss with $\alpha = 1.5$ improved the model's tendency to provide slightly conservative (earlier) RUL predictions, which is preferable for maintenance scheduling.

\section{Conclusion}

This paper presented a comprehensive approach for bearing remaining useful life prediction combining advanced signal processing with deep learning. Our methodology integrates band-pass filtering, Hilbert transform-based envelope analysis, and Db4 wavelet decomposition to extract physics-informed features from multi-channel vibration signals. The CNN-LSTM architecture effectively learns both spatial and temporal patterns from these features, enabling accurate RUL prediction.

Validated on the KSPHM-KIMM 2025 Bearing RUL Prediction Challenge, our approach achieved 7th place among 70 participating teams, demonstrating its effectiveness for real-world bearing prognosis applications. The key contributions of this work include the systematic integration of domain knowledge through signal processing with data-driven deep learning, and the identification of optimal feature combinations (D4\_RMS, D5\_RMS, D5\_Entropy, envelope RMS) for bearing degradation monitoring.

Future work will focus on: (1) incorporating attention mechanisms to improve interpretability, (2) extending the approach to handle variable operating conditions, and (3) developing transfer learning strategies for adapting the model to new bearing types with limited training data.

\section*{Acknowledgment}

The authors would like to thank the Korea Society for Prognostics and Health Management (KSPHM) and the Korea Institute of Machinery and Materials (KIMM) for organizing the 2025 Bearing RUL Prediction Challenge and providing the dataset.

\bibliographystyle{IEEEtran}
\bibliography{references}

\end{document}
